\documentclass[letterpaper, 10 pt, conference]{ieeeconf}  % Comment this line out
                                                          % if you need a4paper
%\documentclass[a4paper, 10pt, conference]{ieeeconf}      % Use this line for a4
                                                          % paper

\IEEEoverridecommandlockouts                              % This command is only
                                                          % needed if you want to
                                                          % use the \thanks command
\overrideIEEEmargins

\usepackage{graphicx}
\usepackage{lipsum}  
\usepackage{xcolor}
\usepackage[hidelinks]{hyperref}

\newcommand{\gepadeu}{{\sc GePaDeU}}


\graphicspath{ {images/} }

\title{\LARGE \bf
\textsc{GePaDeU} --  A Multi-layer Corpus of {\sc Ge}rman {\sc Pa}rliamentary {\sc De}bates with {\sc U}nified Semantic and Pragmatic Annotations 
}

\begin{document}


% https://codingdisciple.com/bootstrap-hypothesis-testing.html

\maketitle
\thispagestyle{empty}
\pagestyle{empty}

%%%%%%%%%%%%%%%%%%%%%%%%%%%%%%%%%%%%%%%%%%%%%%%%%%%%%%%%%%%%%%%%%%%%%%%%%%%%%%%%
\section{Motivation For Dataset Creation}

\textcolor{blue}{\subsection{Why was the dataset created?}} % (e.g., was there a specific task in mind? was there a specific gap that needed to be filled?)}}

The corpus has been created for corpus- and discourse-linguistic studies
of political communication as well as for research
in the political and social sciences.

The main purpose of the annotations is to serve as training data for supervised classification or to benchmark language models for the different semantic and pragmatic extraction and classification tasks.



\textcolor{blue}{\subsection{Has the dataset been used already? If so, where are the results so others can compare
(e.g., links to published papers)?}}

The raw data included in \gepadeu{} is freely availabe and has already been used in different projects and publications, mostly in the field of political science.


\textcolor{blue}{\subsection{What (other) tasks could the dataset be used for?}}

We expect that \gepadeu{} will be interesting for corpus-based investigations of political communication as well as for discourse analyses.


\textcolor{blue}{\subsection{Who funded the creation dataset?}}

The research has been funded by  (anonymous/will be added after the review process).
%Ministry of Science, Research and the Arts Baden Württemberg, Germany.
%  funded by the German Research Foundation (DFG) under the UNCOVER project (PO1900/7-1 and RE3536/3-1).
%Diese Forschungsarbeit von (Autorin/Autor) wurde durch das Ministerium für Wissenschaft, Forschung und Kunst Baden-Württemberg unterstützt.

  

\vspace{1.8cm}
%\pagebreak

\section{Dataset Composition}

\textcolor{blue}{\subsection{What are the instances?}}
%(that is, examples; e.g., documents, images, people, countries) Are there multiple types of instances? (e.g., movies, users, ratings; people, interactions between them; nodes, edges)}}

\gepadeu{} is a text corpus that includes political speeches by members of the German parliament. We provide the data in a json format.
Each json file includes a single speech with annotations for the different layers. 


\textcolor{blue}{\subsection{How many instances are there in total?}}% (of each type, if appropriate)?}}

The dataset includes text from 267 speeches held in the German Bundestag by 196 different speakers (213,617 tokens). The time frame covers the 19th legislative term (2017--2021).
%The data has been split into train, development and test data and has been converted into a ``flat'' version where nested annotations have been ignored (e.g., for nested coordination ``[[children $_{\sc pAge}$] and [adolescents $_{\sc pAge}$] $_{\sc pAge}$]'', we only consider the largest span ``[children and adolescents $_{\sc pAge}$]''). This version of \gepadeu{} includes 7,422 annotated mentions (22,479 annotated tokens). 

%We will also release a version of \gepadeu{} that includes all nested annotations.


% \textcolor{blue}{\subsection{What data does each instance consist of?}}
% %“Raw” data (e.g., unprocessed text or images)? Features/attributes? Is there a label/target associated with instances? If the instances related to people, are subpopulations identified (e.g., by age, gender, etc.) and what is their distribution?}}
% 
% Each instance consists of the text of one paragraph in a tabular format, with one token per line (similar to the CoNLL format for Named Entity Recognition (NER) data.
% Each line includes a paragraph and token id, the word form, the annotations for level 1-3 (our annotation scheme includes hierarchical annotations on three levels) and meta-information (file name, speaker name, party affiliation for the speaker, date of the debate and speech id).
 


\textcolor{blue}{\subsection{Is there a label or target associated with each instance? If so, please
provide a description.}}


For more information on the different annotation layers, please refer to the annotation guidelines in the supplementary materials and our paper submission.


\textcolor{blue}{\subsection{Is any information missing from individual instances?}}
%If so, please provide a description, explaining why this information is missing (e.g., because it was unavailable). This does not include intentionally removed information, but might include, e.g., redacted text.}}

The dataset was created from the transcripts of the parliamentary debates and should thus be considered as a normalised version of the original speech data. We do not include the audio files in the corpus (those are, however, available at \url{https://www.bundestag.de/dokumente/textarchiv/}. 
We also removed all comments and interjections from the speeches. 



\textcolor{blue}{\subsection{Are relationships between individual instances made explicit?}} % (e.g., users’ movie ratings, social network links)? If so, please describe how these relationships are made explicit.}}

The relation between individual speakers can be inferred through the meta-information provided in the data (e.g., party affiliation of the different speakers). The information on date and agenda item also allows to reconstruct which speeches have been given on the same day and topic (however, the topic itself is only specified on an abstract leve, e.g., ``agenda item 1'').


\textcolor{blue}{\subsection{Does the dataset contain all possible instances or is it a sample (not necessarily random) of instances from a larger set?}}
\label{sub:sampling}
%If the dataset is a sample, then what is the larger set? Is the sample representative of the larger set (e.g., geographic coverage)? If so, please describe how this representativeness was validated/verified. If it is not representative of the larger set, please describe why not (e.g., to cover a more diverse range of instances, because instances were withheld or unavailable).}}

The dataset is a sample of the speeches from the 19th legislative term (2017--2021) of the German Bundestag. 
The distribution of topics in \gepadeu{} is not representative of the larger data but has been sampled to cover a more diverse range of topics, with contributions from all parties distributed over the whole legislative term. Below, we describe the sampling procedure in more detail.
 
\paragraph{Sampling procedure}
We extracted a dataset of parliamentary debates from the German Bundestag, covering a time period from the 19th legislative term (2017 to 2021).\footnote{The data is freely available from \url{https://www.bundestag.de/services/opendata}, the Open Data service of the German Bundestag.}
The corpus includes speeches by 807 different speakers, with over 900,000
sentences and over 16 mio tokens. 
From this corpus, we selected individual speeches for annotation as follows.
Our goal was to create a gold standard, controlled for topic and including speeches for each of the political parties. In addition, we wanted the texts to be evenly distributed over the time span of the legislative term (2017--2021). 
To achieve this goal, we selected specific agenda items that covered a range of topics, 
and then sampled all speeches that belong to this specific agenda item, to increase the comparability of the contributions made by speakers from different parties.

\paragraph{CAP topics}
We based our topic selection on the coding scheme developed in the Comparative Agendas Project (CAP) \cite{bevan:2019}. 
The coding scheme includes 21 major topics (see Table \ref{tab:cap}) and more than 200 fine-grained subtopics. The topics we selected have been annotated as major CAP topics, which allowed us to use the annotated CAP data to train a topic classifier.

\begin{table}[t]
\begin{center}\small
\begin{tabular}{ll}
\hline
1 & Cultural Policy Issues \\
2 & Defense \\
3 & Domestic Macroeconomic Issues \\
4 & Education \\
5 & Environment \\
6 & Health  \\
7 & Immigration and Refugee Issues \\
8 & Law, Crime, and Family Issues  \\
\hline
\end{tabular}
\end{center}
\caption{Major topics from the Comparative Agendas Project that we sampled to be included in our data set.}\label{tab:cap} 
\end{table}

\paragraph{Training a CAP topic classifier}
For training data, we used the Parliamentary Question Database\footnote{\url{https://www.comparativeagendas.net/datasets_codebooks}}, a data set with more than 10,000 major and minor interpellations posed by parliamentarians to the government. 
The data set ranges over the 8th to the 15th legislative periods (1976--2005). Each interpellation has been assigned to a major and a minor topic, according to the CAP coding scheme.

Before training, we did some standard preprocessing and clean-up of the data where we lower-cased the text and used a number of regular expressions to remove non-ascii characters, listings of politicians' names, header and footer information and so on. We also removed stopwords and punctuation and extracted a tokenised and lemmatised version of the speeches.\footnote{For lemmatisation, we used the spaCy library: \url{https://spacy.io} with the \texttt{de\_core\_news\_sm} model.} 
This resulted in a training set with 10,033 interpellations, with an average length of 388 tokens per interpellation. 
We then trained a feature-based classifier, based on tf-idf weighted bag-of-words (BOW) features. We experimented with different classifiers provided by the scikit-learn library\footnote{\url{https://scikit-learn.org}} and found that the linear SVM gave us best results for predicting topics on the interpellations. 
For the 21 major topics, our classifier achieves a micro F1 of 72.9\% on the indomain interpellation data.

\paragraph{Sampling based on predicted CAP topics} We then used the classifier to predict topics for each speech in the parliamentary debates, after applying the same preprocessing steps to the data. This gives us topic predictions for each individual speech.
To guide our sampling process, we aggregated the predictions for all speeches belonging to the same agenda item. We call the topic based on a ``majority vote'' for each agenda item the {\em major topic} of the agenda. 
Our assumption is that all speeches given on the same agenda item should belong to the same major topic. 
As a result, we obtained a distribution of topics over all speeches for each respective agenda item. 
We sorted the predictions and {\em manually selected and validated} agenda items for each of the CAP topics in Table \ref{tab:cap}, where the majority of the speeches for this agenda item have been predicted as belonging to this topic. 

We only selected agenda items where each of the political parties participated in the debate, and also aimed at selecting items that are roughly evenly distributed over the time period of the legislative term, to {\em ensure that our data set is as representative as possible, covering a range of different topics, distributed over the whole legislative term and including speeches from all different parties on the same set of topics}.





\textcolor{blue}{\subsection{Are there recommended data splits (e.g., training, development/validation, testing)?}} % If so, please provide a description of these splits, explaining the rationale behind them.}}

The release of the unified corpus does not include data splits for training, development and testing. 
For the splits for the different prediction tasks, we refer users to the repositories for the individual annotation layers (links will be added after the review process).




\textcolor{blue}{\subsection{Are there any errors, sources of noise, or redundancies in the dataset?}} % If so, please provide a description.}}

%The annotations will probably include the usual inconsistencies that are typical for human annotation. 
While we removed comments from the speeches to avoid including speech events that have been produced by persons other than the speaker, the speeches might include some interposed questions or closing remarks not properly marked in the XML version of the data.



\textcolor{blue}{\subsection{Is the dataset self-contained, or does it link to or otherwise rely on external resources (e.g., websites, tweets, other datasets)?}} % If it links to or relies on external resources, a) are there guarantees that they will exist, and remain constant, over time; b) are there official archival versions of the complete dataset (i.e., including the external resources as they existed at the time the dataset was created); c) are there any restrictions (e.g., licenses, fees) associated with any of the external resources that might apply to a future user? Please provide descriptions of all external resources and any restrictions associated with them, as well as links or other access points, as appropriate.}}

The dataset is self-contained and does not rely on other external resources. But note that the audio and video data for the speeches can also be accessed at \url{https://www.bundestag.de/}.

The raw data is in the public domain. The annotated version of the data will be made available under the Creative Commons BY-SA 4.0 license (\url{https://creativecommons.org/licenses/by-sa/4.0/}).


\vspace{1.4cm}

\section{Collection Process}
\label{sec:collection}

\textcolor{blue}{\subsection{How was the data collected?}}

The data has been downloaded from the open data service of the German Bundestag who provide the transcripts of all recent debates in XML format: \url{https://www.bundestag.de/services/opendata}.




\textcolor{blue}{\subsection{If the dataset is a sample from a larger set, what was the sampling strategy?}} %(e.g., deterministic, probabilistic with specific sampling probabilities)?}}


See Section \ref{sub:sampling}, G.



\textcolor{blue}{\subsection{Who was involved in the data collection process (e.g., students, crowdworkers, contractors) and how were they compensated (e.g., how much were crowdworkers paid)?}}

The data has been annotated by student assistants and co-authors of the paper. For  more information, please refer to the paper.





\textcolor{blue}{\subsection{Over what timeframe was the data collected?}} % Does this timeframe match the creation timeframe of the data associated with the instances (e.g., recent crawl of old news articles)? If not, please describe the timeframe in which the data associated with the instances was created.}}

The data was collected in January 2021 and the different annotation layers have been completed 
between January 2021 and September 2024.


\vspace{0.8cm}

\section{Data Preprocessing}

\textcolor{blue}{\subsection{Was any preprocessing/cleaning/labeling of the data done?}} 
%(e.g., discretization or bucketing, tokenization, part-of-speech tagging, SIFT feature extraction, removal of instances, processing of missing values)? If so, please provide a description. If not, you may skip the remainder of the questions in this section.}}

We removed comments and tokenized the data, using spacy (\texttt{de\_core\_news\_sm}).


\textcolor{blue}{\subsection{Was the “raw” data saved in addition to the preprocessed/cleaned/labeled data (e.g., to support unanticipated future uses)?}} % If so, please provide a link or other access point to the “raw” data.}}

The raw data is available from \url{https://www.bundestag.de/services/opendata} in XML format.


\textcolor{blue}{\subsection{Is the software used to preprocess/clean/label the instances available?}} % If so, please provide a link or other access point.}}

We used spaCy for tokenization \url{https://spacy.io/}, with the German \texttt{de\_core\_news\_sm} model.

 

 
\vspace{0.8cm}

\section{Dataset Distribution}

\textcolor{blue}{\subsection{How will the dataset be distributed?}} 
 
We will make the data available  via our university's GitHub account:\\ \url{https://github.com/to-be-added}.
%umanlp/mope.git}.



\textcolor{blue}{\subsection{When will the dataset be released/first distributed?}}
%What license (if any) is it distributed under?}}

\gepadeu{} will be released upon publication of our research paper ``GePaDeU – A Multi-layer Corpus of German Parliamentary Debates with Rich Semantic and Pragmatic Annotations''  that introduces the dataset.
The data will be published under the Creative Commons BY-SA 4.0 license (\url{https://creativecommons.org/licenses/by-sa/4.0/}).


\textcolor{blue}{\subsection{Are there any copyrights on the data?}}

No.

\textcolor{blue}{\subsection{Are there any fees or access/export restrictions?}}

No.

\bigskip


\section{Dataset Maintenance}

\textcolor{blue}{\subsection{Who is supporting/hosting/maintaining the dataset?}}

The dataset will be distributed via the GitHub account of [anonymised] university.


\textcolor{blue}{\subsection{Will the dataset be updated?}}

Probably not.




\textcolor{blue}{\subsection{If the dataset becomes obsolete how will this be
communicated?}}

We do not foresee a scenario where the dataset will become obsolete.


\textcolor{blue}{\subsection{Is there a repository to link to any/all papers/systems that use this dataset?}}
 
No.


\textcolor{blue}{\subsection{If others want to extend/augment/build on this
dataset, is there a mechanism for them to do so?}}

%If so, is there a process for tracking/assessing the quality of those contributions. What is the process for communicating/distributing these contributions to users?}}

The data is available via the Creative Commons BY-SA 4.0 license, so others may extend/augment/build on this dataset, given that they also make the new resource available under the same license.


\pagebreak  


\section{Legal and Ethical Considerations}

\textcolor{blue}{\subsection{Were any ethical review processes conducted (e.g., by an institutional review board)?}} % If so, please provide a description of these review processes, including the outcomes, as well as a link or other access point to any supporting documentation.}}

No.

\textcolor{blue}{\subsection{Does the dataset contain data that might be considered confidential?}}

No.

\textcolor{blue}{\subsection{Does the dataset contain data that, if viewed directly, might be offensive, insulting, threatening, or might otherwise cause anxiety?}} % If so, please describe why}}

The data might include racist and discriminating statements by specific politicians that might be considered as offensive to the user.



\textcolor{blue}{\subsection{Does the dataset relate to people?}} % If not, you may skip the remaining questions in this section.}}

Yes.

\textcolor{blue}{\subsection{Does the dataset identify any subpopulations (e.g., by age, gender)?}} %If so, please describe how these subpopulations are identified and provide a description of their respective distributions within the dataset.}}

The dataset includes speeches by members of the German parliament, held in the Bundestag.
The data collection was conducted by the Bundestag itself and all speakers were aware of the data collection and consented to it. 


\textcolor{blue}{\subsection{Is it possible to identify individuals (i.e., one or more natural persons), either directly or indirectly (i.e., in combination with other
data) from the dataset?}} % If so, please describe how.}}

Yes, all speakers are known.


%\textcolor{blue}{\subsection{Does the dataset contain data that might be considered sensitive in any way (e.g., data that reveals racial or ethnic origins, sexual orientations, religious beliefs, political opinions or union memberships, or locations; financial or health data; biometric or genetic data; forms of government identification, such as social security numbers; criminal history)? If so, please provide a description.}}


%\textcolor{blue}{\subsection{Did you collect the data from the individuals in question directly, or obtain it via third parties or other sources (e.g., websites)?}}



\textcolor{blue}{\subsection{Were the individuals in question notified about the data collection?}}
%If so, please describe (or show with screenshots or other information) how notice was provided, and provide a link or other access point to, or otherwise reproduce, the exact language of the notification itself.}}

The data collection was conducted by the German Bundestag and all speakers were aware of the data collection and consented to it. In addition, the recordings of all debates are freely available on the Bundestag website.





%\textcolor{blue}{\subsection{Did the individuals in question consent to the collection and use of their data? If so, please describe (or show with screenshots or other information) how consent was requested and provided, and provide a link or other access point to, or otherwise reproduce, the exact language to which the individuals consented.}}


%\textcolor{blue}{\subsection{If consent was obtained, were the consenting individuals provided with a mechanism to revoke their consent in the future or for certain uses? If so, please provide a description, as well as a link or other access point to the mechanism (if appropriate).}}


\textcolor{blue}{\subsection{Has an analysis of the potential impact of the dataset and its use on data subjects (e.g., a data protection impact analysis) been conducted?}} % If so, please provide a description of this analysis, including the outcomes, as well as a link or other access point to any supporting documentation.}}

No.



\medskip


\bibliographystyle{IEEEtran}
\bibliography{refs}


\end{document}

